\documentclass[aspectratio=169]{beamer}

\usetheme{Madrid}
\usecolortheme{default}

\usepackage{amsmath,amssymb,mathtools}
\usepackage{physics}
\usepackage{bm}

\title{Derivatives of the Variational Energy}
\subtitle{Gradient and Hessian formulas used in Variational Monte Carlo}
\author{Morten Hjorth-Jensen}
\date{Spring 2026}

\newcommand{\avg}[1]{\left\langle #1 \right\rangle}
\newcommand{\psia}{\psi_{\bm{\alpha}}}
\newcommand{\El}{E_L}
\newcommand{\Oa}{O_\alpha}
\newcommand{\Ob}{O_\beta}
\newcommand{\Oab}{O_{\alpha\beta}}

\begin{document}

%========================================================
\begin{frame}
\titlepage
\end{frame}

%========================================================
\begin{frame}{Setup and notation}
Let $\psia(x)$ be a (real-valued) variational wavefunction depending on parameters
$\bm{\alpha}=(\alpha_1,\dots,\alpha_P)$ and $x$ denote many-body coordinates.

\begin{block}{Variational energy (Rayleigh quotient)}
\[
E(\bm{\alpha})=\frac{\bra{\psia}H\ket{\psia}}{\bra{\psia}\psia\rangle}
=
\frac{\int dx\, \psia(x)\, [H\psia](x)}{\int dx\, \psia(x)^2}.
\]
\end{block}

Define normalization $\mathcal{N}(\bm{\alpha}) := \int dx\, \psia(x)^2$.
\end{frame}

%========================================================
\begin{frame}{Local energy and sampling distribution}
\textbf{Local energy:}
\[
\El(x;\bm{\alpha}) := \frac{[H\psia](x)}{\psia(x)}.
\]

\textbf{Sampling distribution:}
\[
p_{\bm{\alpha}}(x) := \frac{\psia(x)^2}{\mathcal{N}(\bm{\alpha})},
\qquad
\int dx\, p_{\bm{\alpha}}(x)=1.
\]

Then
\[
E(\bm{\alpha})=\int dx\, p_{\bm{\alpha}}(x)\, \El(x;\bm{\alpha})
=\avg{\El}.
\]
\end{frame}

%========================================================
\section{First derivative (gradient)}

%========================================================
\begin{frame}{Differentiate the Rayleigh quotient}
Write
\[
E=\frac{A}{B},
\qquad
A=\int dx\, \psia\,H\psia,\quad B=\int dx\,\psia^2.
\]
Then
\[
\partial_\alpha E = \frac{(\partial_\alpha A)B - A(\partial_\alpha B)}{B^2}.
\]

Compute
\[
\partial_\alpha B = \int dx\, 2\psia\,\partial_\alpha \psia,
\]
and using Hermiticity of $H$,
\[
\partial_\alpha A
=
\int dx\, \Big[(\partial_\alpha \psia)(H\psia)+\psia\,H(\partial_\alpha\psia)\Big]
=
2\int dx\, (\partial_\alpha\psia)\,H\psia.
\]
\end{frame}

%========================================================
\begin{frame}{Gradient in expectation-value form}
Insert into the quotient rule and use $A=EB$:
\[
\partial_\alpha E
=
2\left[
\frac{\int dx\, (\partial_\alpha \psia)\,H\psia}{\int dx\,\psia^2}
-
E\,\frac{\int dx\, \psia\,\partial_\alpha\psia}{\int dx\,\psia^2}
\right].
\]

Rewrite with $\El=(H\psia)/\psia$:
\[
\frac{\int dx\, (\partial_\alpha \psia)\,H\psia}{\int dx\,\psia^2}
=
\frac{\int dx\, \psia^2\left(\frac{\partial_\alpha\psia}{\psia}\right)\El}{\int dx\,\psia^2}
=
\avg{\left(\frac{\partial_\alpha\psia}{\psia}\right)\El}.
\]

Thus
\[
\boxed{
\partial_\alpha E
=
2\left(
\avg{\left(\frac{\partial_\alpha\psia}{\psia}\right)\El}
-\avg{\frac{\partial_\alpha\psia}{\psia}}\avg{\El}
\right).
}
\]
\end{frame}

%========================================================
\begin{frame}{Log-derivative identity (important simplification)}
A key algebraic identity is
\[
\boxed{\frac{\partial_\alpha \psia(x)}{\psia(x)} = \partial_\alpha \ln \psia(x).}
\]

\begin{block}{Definition (log-derivative observable)}
\[
\Oa(x):=\partial_\alpha \ln \psia(x).
\]
Then the gradient becomes the covariance form
\[
\boxed{
\partial_\alpha E = 2\left(\avg{\Oa \El}-\avg{\Oa}\avg{\El}\right)
=2\,\mathrm{Cov}(\Oa,\El).
}
\]
\end{block}

\smallskip
\emph{Complex case (brief remark):} if $\psi$ is complex, define $O_\alpha=\partial_\alpha\ln\psi$ and use
\[
\partial_\alpha E = 2\,\Re\!\left(\avg{O_\alpha^\ast \El}-\avg{O_\alpha^\ast}\avg{\El}\right).
\]
\end{frame}

%========================================================
\section{Second derivative (Hessian)}

%========================================================
\begin{frame}{Key identity: derivative of an expectation value}
Let
\[
\avg{A}=\int dx\, p_{\bm{\alpha}}(x)\,A(x;\bm{\alpha}),
\qquad
p_{\bm{\alpha}}(x)=\frac{\psia(x)^2}{\mathcal{N}}.
\]

Then for real $\psia$ one can show
\[
\boxed{
\partial_\beta \avg{A}
=
\avg{\partial_\beta A}
+2\left(\avg{A\,\Ob}-\avg{A}\avg{\Ob}\right),
}
\qquad
\Ob=\partial_\beta \ln \psia.
\]

\begin{itemize}
\item The first term accounts for explicit $\bm{\alpha}$-dependence of $A$.
\item The second term accounts for how the sampling distribution changes with $\bm{\alpha}$.
\end{itemize}
\end{frame}

%========================================================
\begin{frame}{Useful derivative: $\partial_\beta E_L$}
Since $\El=(H\psia)/\psia$, its parameter derivative is
\[
\partial_\beta \El
=
\partial_\beta\left(\frac{H\psia}{\psia}\right)
=
\frac{H(\partial_\beta \psia)}{\psia}
-\El\,\frac{\partial_\beta \psia}{\psia}.
\]

Using $\Ob=\partial_\beta\ln\psia=(\partial_\beta\psia)/\psia$:
\[
\boxed{
\partial_\beta \El
=
\frac{H(\partial_\beta \psia)}{\psia}
-\El\,\Ob.
}
\]

\begin{block}{Notation}
Also define the second log-derivative
\[
\boxed{
\Oab(x):=\partial_\beta \Oa(x)=\partial_\alpha\partial_\beta \ln \psia(x).
}
\]
\end{block}
\end{frame}

%========================================================
\begin{frame}{Hessian: start from the covariance gradient}
We have
\[
\partial_\alpha E = 2\left(\avg{\Oa \El}-\avg{\Oa}\avg{\El}\right).
\]

Differentiate w.r.t.\ $\beta$:
\[
\partial_\beta\partial_\alpha E
=
2\left[
\partial_\beta \avg{\Oa \El}
-\partial_\beta(\avg{\Oa}\avg{\El})
\right].
\]

We now apply the expectation-derivative identity to each expectation value.
\end{frame}

%========================================================
\begin{frame}{Explicit Hessian formula (practical VMC form)}
Using
\(
\partial_\beta\avg{A}=\avg{\partial_\beta A}+2(\avg{A\Ob}-\avg{A}\avg{\Ob}),
\)
with $A=\Oa\El$, $A=\Oa$, and $A=\El$, we obtain:

\[
\boxed{
\begin{aligned}
\partial_\beta\partial_\alpha E
=2\Big[&
\avg{\Oab\,\El}+\avg{\Oa\,\partial_\beta \El}
-\avg{\Oab}\avg{\El}-\avg{\Oa}\avg{\partial_\beta \El} \\
&\quad
+2\Big(
\avg{\Oa\,\El\,\Ob}-\avg{\Oa\,\El}\avg{\Ob}
-\avg{\Oa\,\Ob}\avg{\El}-\avg{\Oa}\avg{\El\,\Ob}
+2\,\avg{\Oa}\avg{\Ob}\avg{\El}
\Big)
\Big].
\end{aligned}
}
\]

\begin{itemize}
\item First line: ``explicit'' derivatives ($\Oab$ and $\partial_\beta \El$).
\item Second line: ``distribution-derivative'' terms (a third-order cumulant-like combination).
\end{itemize}
\end{frame}

%========================================================
\begin{frame}{Compact cumulant interpretation (optional)}
Define centered variables
\[
\delta \Oa = \Oa-\avg{\Oa},\quad
\delta \Ob=\Ob-\avg{\Ob},\quad
\delta \El=\El-\avg{\El}.
\]

Then the gradient is
\[
\partial_\alpha E = 2\,\avg{\delta\Oa\,\delta\El}.
\]

The Hessian can be viewed as
\[
\partial_\beta\partial_\alpha E
=
2\,\mathrm{Cov}(\Oab,\El)
+2\,\mathrm{Cov}(\Oa,\partial_\beta \El)
+4\,\kappa_3(\Oa,\Ob,\El),
\]
where $\kappa_3$ denotes the third cumulant (connected three-point function)
\[
\kappa_3(\Oa,\Ob,\El)=\avg{\delta\Oa\,\delta\Ob\,\delta\El}.
\]

\smallskip
This form highlights why Hessian estimates are typically noisier than gradients in Monte Carlo.
\end{frame}

%========================================================
\begin{frame}{Monte Carlo estimators (implementation hint)}
Given samples $\{x^{(m)}\}_{m=1}^M \sim p_{\bm{\alpha}}$,
compute
\[
\Oa^{(m)}=\partial_\alpha \ln\psia(x^{(m)}),\quad
\Ob^{(m)}=\partial_\beta \ln\psia(x^{(m)}),\quad
\El^{(m)}=\El(x^{(m)}),
\]
and, if needed,
\[
\Oab^{(m)}=\partial_\alpha\partial_\beta \ln\psia(x^{(m)}),\qquad
(\partial_\beta \El)^{(m)}=\left.\partial_\beta \El\right|_{x^{(m)}}.
\]

\begin{block}{Practical remark}
Many VMC optimizers avoid explicit Hessians (noisy) and instead use
stochastic reconfiguration / natural-gradient methods that only require first derivatives.
\end{block}
\end{frame}

%========================================================
\begin{frame}{Summary}
\begin{itemize}
\item Define $p(x)\propto \psi(x)^2$ and $\El=(H\psi)/\psi$ so $E=\avg{\El}$.
\item Log-derivative identity:
\[
\frac{\partial_\alpha\psi}{\psi}=\partial_\alpha\ln\psi.
\]
\item Gradient:
\[
\partial_\alpha E = 2\left(\avg{\Oa\El}-\avg{\Oa}\avg{\El}\right)=2\,\mathrm{Cov}(\Oa,\El).
\]
\item Hessian: obtained by differentiating the covariance and accounting for changes in $p(x)$,
yielding an explicit VMC-ready expression involving $\Oab$ and $\partial_\beta \El$.
\end{itemize}
\end{frame}

\end{document}
